\section{Markov Chain Monte Carlo - MCMC}

\subsection{Leituras Recomendadas}
\begin{frame}{Markov Chain Monte Carlo - MCMC - Leituras Recomendadas}
    \begin{vfilleditems}
        \item \textcite{gelman2013bayesian}
        \begin{vfilleditems}
            \item Capítulo 10: Introduction to Bayesian computation
            \item Capítulo 11: Basics of Markov chain simulation
            \item Capítulo 12: Computationally efficient Markov chain simulation
        \end{vfilleditems}
        \item \textcite{mcelreath2020statistical} - Capítulo 9: Markov Chain Monte Carlo
        \item \textcite{neal2011mcmc}
        \item \textcite{betancourtConceptualIntroductionHamiltonian2017}
        \item \textcite{gelman2020regression} - Capítulo 22, Seção 22.8: Computational efficiency
        \item \textcite{storopoli2021estatisticabayesianaR} - Markov Chain Monte Carlo
        \item \textcite{chibUnderstandingMetropolisHastingsAlgorithm1995}
        \item \textcite{casellaExplainingGibbsSampler1992}
    \end{vfilleditems}
\end{frame}

\begin{frame}{Métodos de Monte Carlo}
  \begin{columns}
    \begin{column}{0.8\textwidth}
      \begin{vfilleditems}
        \item \href{http://mc-stan.org/}{\texttt{Stan}} é uma homenagem ao matemático
        Stanislaw Ulam, que participou do projeto Manhattan e ao tentar calcular o
        processo de difusão de neutrons para a bomba de hidrogênio acabou criando uma
        classe de métodos chamada \textbf{Monte Carlo} \parencite{eckhardtStanUlamJohn1987}.
        \item Métodos de Monte Carlo possuem como conceito subjacente o uso a aleatoriedade
        para resolver problemas que podem ser determinísticos em princípio.
        Eles são freqüentemente usados em problemas físicos e matemáticos e
        são mais úteis quando é difícil ou impossível usar outras abordagens.
      \end{vfilleditems}
    \end{column}
    \begin{column}{0.2\textwidth}
      \centering
      \includegraphics[width=0.9\columnwidth]{stanislaw.jpg}
    \end{column}
  \end{columns}
\end{frame}

\begin{frame}{História dos Métodos de Monte Carlo\footnote{para quem se interessou, a história se encontra em \textcite{eckhardtStanUlamJohn1987}}}
  \begin{columns}
    \begin{column}{0.8\textwidth}
      \begin{vfilleditems}
        \item A ideia do método veio enquanto jogava paciência durante sua
        recuperação de uma cirurgia, Ulam pensou em jogar centenas de jogos para
        estimar estatisticamente a probabilidade de um resultado bem-sucedido
        \item Ulam descreveu a ideia para John von Neumann em 1946
        \item \small Por ser secreto, o trabalho de von Neumann e Ulam exigia um codinome.
        Um colega de von Neumann e Ulam, Nicholas Metropolis, sugeriu usar o nome Monte Carlo,
        que se refere ao Casino Monte Carlo em Mônaco, onde o tio de Ulam (Michał Ulam)
        pedia dinheiro emprestado a parentes para jogar.
      \end{vfilleditems}
    \end{column}
    \begin{column}{0.2\textwidth}
      \centering
      \includegraphics[width=0.9\columnwidth]{stanislaw.jpg}
    \end{column}
  \end{columns}
\end{frame}

\subsection{Por quê Precisamos de MCMC?}
\begin{frame}{Por quê Precisamos de MCMC?}
  A principal barreira computacional para estatística Bayesiana é o denominador
  $P(\text{data})$ da fórmula de Bayes:

  $$P(\theta \mid \text{data})=\frac{P(\theta) \cdot P(\text{data} \mid \theta)}{P(\text{data})}$$

  Em casos discretos podemos fazer o denominador virar a soma de todos os parâmetros
  usando a \textbf{regra da cadeia de probabilidade} (\textit{chain rule}):

  $$P(A,B \mid C)=P(A \mid B,C) \times P(B \mid C)$$

  Isto também é chamado de \textbf{marginalização}:

  $$P(\text{data})=\sum_{\theta} P(\text{data} \mid \theta) \times P(\theta)$$
\end{frame}

\begin{frame}{Por quê Precisamos de MCMC?}
  Porém no caso de valores contínuos o denominador $P(\text{data})$ vira uma integral
  bem grande e complicada de calcular:

  $$P(\text{data})=\int_{\theta} P(\text{data} \mid \theta) \times P(\theta)d \theta$$

  Em muitos casos essa integral vira \textit{intratável} (incalculável) e
  portanto devemos achar outras maneiras de calcular a probabilidade posterior
  $P(\theta \mid \text{data})$ de Bayes sem usar o denominador $P(\text{data})$.
  \vfill
  \Large \textbf{É aqui que entra Métodos de Monte Carlo!}
\end{frame}

\begin{frame}{Para quê serve o denominador $P(\text{data})$}
  Para normalizar a posterior com o intuito de torná-la uma distribuição
  probabilística válida. Isto quer dizer que a soma de todas as probabilidades
  dos eventos possíveis da distribuição devem ser iguais a $1$:
  \begin{vfilleditems}
    \item no caso de distribuição probabilística \textbf{discreta}:
    $$\sum_{\theta} P(\theta \mid \text{data}) = 1$$
    \item no caso de distribuição probabilística \textbf{contínua}:
    $$\int_{\theta} P(\theta \mid \text{data})d \theta = 1$$
  \end{vfilleditems}
\end{frame}

\begin{frame}{Se removermos o denominador de Bayes o que temos?}
  Ao removermos o denominador $(\text{data})$ temos que a posterior
  $P(\theta \mid \text{data})$ é \textbf{proporcional} à \textit{priori}
  multiplicada pela verossimilhança $P(\theta) \cdot P(\text{data} \mid \theta)$:

  $$P(\theta \mid \text{data}) \propto P(\theta) \cdot P(\text{data} \mid \theta)$$

\end{frame}

\subsubsection{Correntes Markov}
\begin{frame}{Método de Montecarlo com Correntes Markov -- (MCMC)}
  Aí que entra \textbf{Método Montecarlo com Correntes Markov}\footnote{do inglês
  \textit{Markov Chain Monte Carlo} (MCMC)}.
  \vfill
  MCMC é uma classe ampla de ferramentas computacionais para aproximação
  de integrais e geração de amostras de uma probabilidade posterior
  \parencite{brooksHandbookMarkovChain2011}.
  \vfill
  MCMC é usada quando não é possível coletar amostras de $\boldsymbol{\theta}$
  direto da distribuição probabilística posterior
  $P(\boldsymbol{\theta} \mid \text{data})$.
  Ao invés disso, nos coletamos amostras de maneira iterativa que a cada passo do
  processo nós esperamos que a distribuição da qual amostramos
  $P^*(\boldsymbol{\theta}^{(*)} \mid \text{data})$
  se torna cada vez mais similar à posterior $P(\boldsymbol{\theta} \mid \text{data})$.
  \vfill
  Tudo isso é para \textbf{eliminar o cálculo} (muitas vezes impossível) do \textbf{denominador} $P(\text{data})$.
\end{frame}

\begin{frame}{Correntes Markov}
  \begin{columns}
    \begin{column}{0.8\textwidth}
      \begin{vfilleditems}
        \item A ideia é \textbf{definir uma corrente Markov ergódica}
        (quer dizer que há uma distribuição estacionária única)
        dos quais o conjunto de estados possíveis é o espaço amostral e a
        distribuição estacionária é a distribuição a ser \textit{aproximada} (ou \textit{amostrada}).
        \item Seja $X_0, X_1, \dots, X_n$ uma simulação da corrente.
        A corrente Markov \textbf{converge à distribuição estacionária de qualquer
        estado inicial} $X_0$ após um \textbf{número suficiente grande de iterações} $r$,
        a distribuição do estado $X_r$ estará similar à distribuição estacionária,
        então podemos usá-la com amostra.
      \end{vfilleditems}
    \end{column}
    \begin{column}{0.2\textwidth}
      \centering
      \includegraphics[width=0.9\columnwidth]{andrei_markov.jpg}
    \end{column}
  \end{columns}
\end{frame}

\begin{frame}{Correntes Markov}
  \begin{columns}
    \begin{column}{0.8\textwidth}
      \begin{vfilleditems}
        \item As correntes Markov possuem uma propriedade que a distribuição probabilística
        do próximo estado depende \textbf{apenas do estado atual e não na sequência
        de eventos que precederam}:
        $P(X_{n+1}=x \mid X_{0},X_{1},X_{2},\ldots ,X_{n}) = P(X_{n+1}=x \mid X_{n})$.
        Essa propriedade é chamada de \textbf{Markoviana}.
        \item Similarmente, repetindo esse argumento com $X_r$ como o ponto inicial,
        podemos usar $X_{2r}$ como amostra, e assim por diante.
        Podemos então usar a sequência de estados $X_r, X_{2r}, X_{3r}, \dots$
        como quase \textbf{amostras independentes} da distribuição estacionária da
        corrente Markov.
      \end{vfilleditems}
    \end{column}
    \begin{column}{0.2\textwidth}
      \centering
      \includegraphics[width=0.9\columnwidth]{andrei_markov.jpg}
    \end{column}
  \end{columns}
\end{frame}

% Idea taken from http://steventhornton.ca/blog/markov-chains-in-latex.html
\begin{frame}{Exemplo de Corrente Markov}
  \centering
  \begin{tikzpicture}
    % Add the states
    \node[state,
    text=yellow,
    minimum size=2cm,
    thick
    ]
    (s) {Sol};
    \node[state,
    right=3cm of s,
    text=blue!30!white,
    minimum size=2cm,
    thick
    ]
    (r) {Chuva};

% Connect the states with arrows
\draw[every loop,
    auto=right,
    line width=1mm,
    >=latex]
  (s) edge[bend right, auto=left]  node {0.6} (r)
  (r) edge[bend right, auto=right] node {0.7} (s)
  (s) edge[loop above]             node {0.4} (s)
  (r) edge[loop above]             node {0.3} (r);
\end{tikzpicture}
\end{frame}

\begin{frame}{Correntes Markov}
  A eficácia dessa abordagem depende em:

  \begin{vfilleditems}
    \item \textbf{o quão grande $r$ deve ser} para garantir uma \textbf{amostra adequadamente boa}
    \item \textbf{poder computacional} requerido para cada iteração da corrente Markov.
  \end{vfilleditems}

  \vfill
  \footnotesize
  Além disso, é costumeiro descartarmos as primeiras iterações do algoritmo pois
  elas costumam não ser representativas da distribuição a ser aproximada.
  Nas iterações iniciais de algoritmos MCMC geralmente a corrente Markov
  está em um processo de aquecimento\footnote{Algumas referências chamam esse processo de \textit{burnin}}
  (\textit{warm-up}) e seu estado está bem distante do ideal para começarmos uma amostragem
  fidedigna.
  \vfill
  Geralmente, recomenda-se que se descarte metade das iterações \parencite{gelmanBasicsMarkovChain2013}.
\end{frame}

\subsection{Algoritmos de MCMC}
\begin{frame}{Algoritmos de MCMC}
  Temos \textbf{MUITOS} algoritmos de MCMC\footnote{Veja a \href{https://en.wikipedia.org/wiki/Markov_chain_Monte_Carlo}{página da Wikipedia para uma listagem completa}}
  Mas aqui vamos cobrir duas classes de algoritmos MCMC:
  \begin{vfilleditems}
    \item Metropolis-Hastings \parencite{metropolisEquationStateCalculations1953, hastingsMonteCarloSampling1970}

    \item Hamiltonian Monte Carlo\footnote{às vezes chamado de \textit{Hybrid Monte Carlo}, especialmente na literatura de Física} \parencite{neal2011mcmc, betancourtConceptualIntroductionHamiltonian2017}
  \end{vfilleditems}
\end{frame}

\begin{frame}{Classe de Algoritmos MCMC -- Metropolis-Hastings}
  Os primeiros algoritmos de MCMC. Usam uma regra de aceitação/rejeição das
  propostas. Caracterizados por propostas oriundas de um passeio aleatório\footnote{\textit{random walk}}
  no espaço amostral. O algoritmo de \textbf{Gibbs} pode ser visto como um
  \textbf{caso especial} do algoritmo de MH porque
  todas as propostas são aceitas \parencite{gelmanIterativeNonIterativeSimulation1992}
  \vfill
  Assintoticamente, possuem uma taxa de aceitação de 23.4\% e o custo de cada iteração é
  $\mathcal{O}(d)$, na qual $d$ é a dimensão do espaço amostral \parencite{beskosOptimalTuningHybrid2013}.
\end{frame}

\begin{frame}{Classe de Algortimos MCMC -- Hamiltonian Monte Carlo}
  Os algoritmos MCMC mais eficientes na atualidade. Tenta evitar o comportamento
  de passeio aleatório introduzindo um vetor de momento auxiliar e
  implementando dinâmicas Hamiltonianas. As propostas são "guiadas"
  para regiões de maior densidade do espaço amostral. Isso faz com que HMC seja
  \textbf{ordens de magnitude mais eficiente que MH e Gibbs}.
  \vfill
  Assintoticamente, possuem uma taxa de aceitação de 65.1\% e o custo de cada iteração é
  $\mathcal{O}(d^{\frac{1}{4}})$, na qual $d$ é a dimensão do espaço amostral \parencite{beskosOptimalTuningHybrid2013}.
\end{frame}

\subsubsection{Metropolis}
% \includegraphics[width=0.9\columnwidth]{nicholas_metropolis.png}
\begin{frame}[fragile]{Algoritmo de Metropolis\footnote{\textcite{metropolisEquationStateCalculations1953}}}
    \SetAlCapFnt{\footnotesize}
    \SetAlCapNameFnt{\footnotesize}
    \begin{algorithm}[H]
    \DontPrintSemicolon
    \SetAlgoNoEnd
    \SetAlgoLined
    Defina um ponto inicial $\boldsymbol{\theta}^{(0)} \in \mathbb{R}^p$ do qual $P\left(\boldsymbol{\theta}^{(0)} \mid \boldsymbol{y} \right) > 0$\;
     \Para{$t = 1, 2, \dots$}{
      Amostra uma proposta $\boldsymbol{\theta}^{(*)}$ de uma distribuição de propostas no tempo $t$, $J_t \left(\boldsymbol{\theta}^{(*)} \mid \boldsymbol{\theta}^{(t-1)} \right)$\;
      Como regra de aceitação/rejeição calcule a proporção das probabilidades:
      $r = \frac{P\left(\boldsymbol{\theta}^{(*)}  \mid \boldsymbol{y} \right)}{P\left(\boldsymbol{\theta}^{(t-1)} \mid \boldsymbol{y} \right)}$\;
      Designe:
      $
        \boldsymbol{\theta}^{(t)} =
          \begin{cases}
          \boldsymbol{\theta}^{(*)} & \text{com probabilidade $\min(r,1)$}\\
          \boldsymbol{\theta}^{(t-1)} & \text{caso contrário}
        \end{cases}
      $\;
     }
     \caption{Metropolis}
    \end{algorithm}
\end{frame}

\begin{frame}{Intuição Visual de Metropolis}
\centering
    \begin{tikzpicture}
        \begin{axis}[every axis plot, line width=2pt,
            ylabel=PDF,
            domain=-4:4,samples=200,
            ymax = 0.6, ytick={0, 0.2, 0.4},
            axis x line*=bottom, % no box around the plot, only x and y axis
            axis y line*=left, % the * suppresses the arrow tips
            enlargelimits=true,
            ] % extend the axes a bit

            \addplot [blue] {gaussian(0, 1)};
            \node[inner sep=0pt] (hikerlower) at (-2,0.13){\Strichmaxerl[2pt]};
            \node[inner sep=0pt] (hikerupper) at (0,0.5){\Strichmaxerl[2pt]};
            \node[inner sep=0pt] (hikerlower2) at (2,0.13){\Strichmaxerl[2pt]};
            \draw[->, red, line width=2pt] (hikerlower) to [out=90,in=135] node[above left] {\large$P=1$} (hikerupper);
            \draw[->, yellow, line width=2pt] (hikerupper) to [out=45,in=135] node[right] {\large$P\approx\frac{0.1}{0.4}\approx\frac{1}{4}$} (hikerlower2);
        \end{axis}
        \end{tikzpicture}

\end{frame}

\subsubsection{Metropolis-Hastings}
\begin{frame}[fragile]{Algoritmo de Metropolis - Hastings\footnote{\textcite{hastingsMonteCarloSampling1970}}}
    \SetAlCapFnt{\footnotesize}
    \SetAlCapNameFnt{\footnotesize}
    \begin{algorithm}[H]
    \DontPrintSemicolon
    \SetAlgoNoEnd
    \SetAlgoLined
    Defina um ponto inicial $\boldsymbol{\theta}^{(0)} \in \mathbb{R}^p$ do qual $P\left(\boldsymbol{\theta}^{(0)} \mid \boldsymbol{y} \right) > 0$\;
     \Para{$t = 1, 2, \dots$}{
      Amostra uma proposta $\boldsymbol{\theta}^{(*)}$ de uma distribuição de propostas no tempo $t$, $J_t \left(\boldsymbol{\theta}^{(*)} \mid \boldsymbol{\theta}^{(t-1)} \right)$\;
      Como regra de aceitação/rejeição calcule a proporção das probabilidades:
      $r = \frac{\frac{P \left(\boldsymbol{\theta}^{(*)} \mid \boldsymbol{y} \right)}{J_t \left(\boldsymbol{\theta}^{(*)} \mid \boldsymbol{\theta}^{(t-1)} \right)}}{\frac{P \left(\boldsymbol{\theta}^{(t-1)} \mid \boldsymbol{y} \right)}{J_t \left(\boldsymbol{\theta}^{(t-1)} \mid \boldsymbol{\theta}^{(*)} \right)}}$\;
      Designe:
      $
        \boldsymbol{\theta}^{(t)} =
          \begin{cases}
          \boldsymbol{\theta}^{(*)} & \text{com probabilidade $\min(r,1)$}\\
          \boldsymbol{\theta}^{(t-1)} & \text{caso contrário}
        \end{cases}
      $\;
     }
     \caption{Metropolis-Hastings}
    \end{algorithm}
\end{frame}

\begin{frame}{Animação Metropolis\footnote{veja Metropolis em ação no \href{https://chi-feng.github.io/mcmc-demo/app.html?algorithm=RandomWalkMH&target=banana}{\texttt{chi-feng/mcmc-demo}}}}
  \centering
  \includemedia[
    width=\linewidth,
    height=0.3\linewidth,
    addresource=rwmh.m4v,
    transparent,
    activate=pageopen,
    passcontext,  %show VPlayer's right-click menu
    flashvars={
      source=rwmh.m4v
      &loop=true
      &scaleMode=stretch
    }
  ]{Animação Metropolis}{http://mirrors.ctan.org/macros/latex/contrib/media9/players/VPlayer.swf}
\end{frame}

\subsubsection{As Limitações de Metropolis e Metropolis-Hastings}

\subsubsection{Gibbs}
\begin{frame}[fragile]{Algoritmo de Gibbs\footnote{\textcite{gemanStochasticRelaxationGibbs1984}}}
    \SetAlCapFnt{\footnotesize}
    \SetAlCapNameFnt{\footnotesize}
    \begin{algorithm}[H]
    \DontPrintSemicolon
    \SetAlgoNoEnd
    \SetAlgoLined
    Defina um ponto inicial $\boldsymbol{\theta}^{(0)} \in \mathbb{R}^p$ do qual $P\left(\boldsymbol{\theta}^{(0)} \mid \boldsymbol{y} \right) > 0$\;
     \Para{$t = 1, 2, \dots$}{
      Designe:
      $ \boldsymbol{\theta}^{(t)} =
        \begin{cases}
        \theta^{(t)}_1 &\sim P \left(\theta_1 \mid \theta^{(0)}_2, \dots, \theta^{(0)}_p \right) \\
        \theta^{(t)}_2 &\sim P \left(\theta_2 \mid \theta^{(t-1)}_1, \dots, \theta^{(0)}_p \right) \\
        &\vdots \\
        \theta^{(t)}_p &\sim P \left(\theta_p \mid \theta^{(t-1)}_1, \dots, \theta^{(t-1)}_{p-1} \right)
     \end{cases}
      $\;
     }
     \caption{Gibbs}
    \end{algorithm}
\end{frame}

\begin{frame}{Animação Gibbs\footnote{Veja Gibbs em ação no \href{https://chi-feng.github.io/mcmc-demo/app.html?algorithm=GibbsSampling&target=banana}{\texttt{chi-feng/mcmc-demo}}}}
  \centering
  \includemedia[
    width=\linewidth,
    height=0.3\linewidth,
    addresource=gibbs.m4v,
    transparent,
    activate=pageopen,
    passcontext,  %show VPlayer's right-click menu
    flashvars={
      source=gibbs.m4v
      &loop=true
      &scaleMode=stretch
    }
  ]{Animação Gibbs}{http://mirrors.ctan.org/macros/latex/contrib/media9/players/VPlayer.swf}
\end{frame}

\subsubsection{As Limitações de Gibbs}

\subsubsection{Hamiltonian Monte Carlo (HMC)}
% Störmer–Verlet or leapfrog
\begin{frame}[fragile]{Algoritmo de HMC\footnote{\textcite{neal2011mcmc}}}
    \SetAlCapFnt{\footnotesize}
    \SetAlCapNameFnt{\footnotesize}
    \begin{algorithm}[H]
    \DontPrintSemicolon
    \SetAlgoNoEnd
    \SetAlgoLined
    Defina um ponto inicial $\boldsymbol{\theta}^{(0)} \in \mathbb{R}^p$ do qual $P\left(\boldsymbol{\theta}^{(0)} \mid \boldsymbol{y} \right) > 0$\;
    Amostre $\boldsymbol{\phi}$ de uma $\text{Normal}(\mathbf{0},\mathbf{M})$\;
    Simultaneamente amostre $\boldsymbol{\theta}$ e $\boldsymbol{\phi}$ com $L$ \textit{leapfrog steps} cada um reduzido por um fator $\epsilon$.\;
    \Para{$1, 2, \dots, L$}{
     Use o gradiente do $\log$ da posterior de $\boldsymbol{\theta}$ para produzir um \textit{half-step} de $\boldsymbol{\phi}$:
     $\boldsymbol{\phi} \leftarrow \boldsymbol{\phi} + \frac{1}{2} \epsilon \frac{d \log P(\boldsymbol{\theta} \mid \boldsymbol{y})}{d \theta}$\;
     Use $\boldsymbol{\phi}$ para atualizar $\boldsymbol{\theta}$:
     $\boldsymbol{\theta} \leftarrow \boldsymbol{\theta} + \epsilon \mathbf{M}^{-1} \boldsymbol{\phi}$\;
     Novamente use o gradiente de $\boldsymbol{\theta}$ para produzir um \textit{half-step} de $\boldsymbol{\phi}$:
     $\boldsymbol{\phi} \leftarrow \boldsymbol{\phi} + \frac{1}{2} \epsilon \frac{d \log P(\boldsymbol{\theta} \mid \boldsymbol{y})}{d \theta}$\;
    }
    Como regra de aceitação/rejeição calcule:
    $r = \frac{P \left(\boldsymbol{\theta}^{(*)} \mid \boldsymbol{y} \right) P \left(\boldsymbol{\phi}^{(*)} \right)}{P \left(\boldsymbol{\theta}^{(t-1)} \mid \boldsymbol{y} \right) P \left(\boldsymbol{\phi}^{(t-1)} \right)}$\;
    Designe:
      $
        \boldsymbol{\theta}^{(t)} =
          \begin{cases}
          \boldsymbol{\theta}^{(*)} & \text{com probabilidade $\min(r,1)$}\\
          \boldsymbol{\theta}^{(t-1)} & \text{caso contrário}
        \end{cases}
      $\;
    \caption{Hamiltonian Monte Carlo (HMC)}
    \end{algorithm}
\end{frame}

\begin{frame}{Animação HMC\footnote{Veja HMC em ação no \href{https://chi-feng.github.io/mcmc-demo/app.html?algorithm=HamiltonianHMC&target=banana}{\texttt{chi-feng/mcmc-demo}}}}
  \includemedia[
    width=\linewidth,
    height=0.3\linewidth,
    addresource=hmc.m4v,
    transparent,
    noplaybutton,
    activate=pageopen,
    passcontext,  %show VPlayer's right-click menu
    flashvars={
      source=hmc.m4v
      &loop=true
      &scaleMode=stretch
    }
  ]{Animação HMC}{http://mirrors.ctan.org/macros/latex/contrib/media9/players/VPlayer.swf}
\end{frame}

\subsubsection{No-U-Turn-Sampler (NUTS)}

\begin{frame}{Algoritmo de NUTS\footnote{\textcite{hoffman2014no}}}
    \SetAlCapFnt{\footnotesize}
    \SetAlCapNameFnt{\footnotesize}
    \begin{algorithm}[H]
    \DontPrintSemicolon
    \SetAlgoNoEnd
    \SetAlgoLined
    \caption{No-U-Turn-Sampler (NUTS)}
    \end{algorithm}
\end{frame}

\begin{frame}{Animação NUTS\footnote{Veja NUTS em ação no \href{https://chi-feng.github.io/mcmc-demo/app.html?algorithm=EfficientNUTS&target=banana}{\texttt{chi-feng/mcmc-demo}}}}
  \centering
  \includemedia[
    width=\linewidth,
    height=0.3\linewidth,
    addresource=nuts.m4v,
    transparent,
    activate=pageopen,
    passcontext,  %show VPlayer's right-click menu
    flashvars={
      source=nuts.m4v
      &loop=true
      &scaleMode=stretch
    }
  ]{Animação NUTS}{http://mirrors.ctan.org/macros/latex/contrib/media9/players/VPlayer.swf}
\end{frame}

\subsubsection{As Limitações de HMC e NUTS}

% Funnel stuff from here https://betanalpha.github.io/assets/case_studies/hierarchical_modeling.html#33_Funnel_Degeneracies_and_Bayesian_Computation
% https://github.com/betanalpha/knitr_case_studies/tree/master/hierarchical_modeling/figures/funnel

\subsection{Convergência de Correntes Markov}

\subsubsection{Métricas de Convergência}

\subsubsection{Visualizações de Convergência}

\subsubsection{O que fazer se Correntes Markov não Convergirem}
